\section{Model Views, Runtime Inspector, and Human Observability}
\label{sec:views}

\subsection{Why topology is derived}

A module graph is an excellent abstraction for reconfiguration algorithms, but it
cannot be canonical physical state. It omits links, joint values, connector frames,
contact, pose, and transient docking lifecycle. Making it independently mutable would
create the duplicated-truth problem \modsim{} is designed to avoid.

Model views are therefore strict, immutable, JSON-safe Pydantic data-transfer objects.
They contain a schema version, source stamp, simulation time, pack identity, and typed
content. They never contain Qt objects, \mujoco{} handles, NumPy arrays, NetworkX
graphs, or references to mutable \code{WorldState}. A client may convert the DTO into
its chosen rendering or analysis library.

\subsection{Factory and builder abstraction}

\code{ModelViewBuilder} is the extension point. A builder declares a stable ID, name,
result type, supported modes, whether it requires a live world, a dependency token,
and a build method. \code{ModelViewFactory} is per-client rather than global. It
registers built-ins, rejects duplicate IDs, lists descriptors, resolves pack recipes,
builds one or all views, and maintains a bounded thread-safe cache.

Robot Pack recipes contain only builder ID, modes, default flag, display name, and
bounded JSON configuration. They contain no executable import path. An unknown builder
remains portable valid data but cannot be constructed until an application registers an
implementation.

\subsection{Module topology graph}

The first built-in view is a runtime multigraph:

\begin{itemize}
  \item every module instance is a node, including isolated modules;
  \item every committed connection is one undirected edge;
  \item parallel connections remain separate edges;
  \item nodes copy module pose and assembly identity;
  \item edges copy endpoints, connection ID, orientation index, and commit time; and
  \item deterministic sorting makes serialization stable.
\end{itemize}

Only \code{DockCommitted} and \code{UndockCommitted} change graph topology. Candidates,
failed attempts, and planned edges are not drawn as committed connections.

\subsection{Revision-aware caching}

The topology builder observes the backend sample sequence because node poses can move,
and the topology revision because edges can change. A cache entry is keyed by source
identity, recipe fingerprint, and builder token. Unrelated event or docking revisions
can update a source stamp without requiring a full graph reconstruction. The cache is
least-recently-used and bounded, retaining the newest views rather than growing with
simulation length.

\subsection{Runtime Inspector frame}

After stepping, the runtime owner builds a \code{RuntimeInspectorFrame} containing:

\begin{itemize}
  \item the immutable graph view;
  \item scenario phase, plan/action index, and human-readable detail;
  \item current runtime metrics;
  \item a contiguous event delta since the last successfully published frame; and
  \item backend, session, time, and source revision information.
\end{itemize}

Graph frames may be refresh-limited, but event deltas are lossless and contiguous. The
event cursor advances only after a frame is successfully constructed for publication.

\subsection{Two execution lanes}

The Runtime Inspector supports two ownership patterns that share the same frame
contract.

\paragraph{Semantic-only lane.}
With \code{--no-viewer}, a dedicated QThread worker owns the runtime, scenario,
backend, and model-view factory. It paces against simulated time and emits immutable
frames to the Qt thread at a configured wall-clock publication rate.

\paragraph{Native-viewer lane.}
With \mujoco{} viewer enabled, a companion process owns the one authoritative runtime
and native viewer. This is required on macOS because Qt and the passive \mujoco{} viewer
cannot both own the Cocoa main GUI thread. The Qt parent displays topology/events; the
child advances physics and displays geometry. They are not two synchronized
simulations.

\begin{figure}[htbp]
\centering
\begin{tikzpicture}[
  node distance=7mm and 18mm,
  proc/.style={draw,rounded corners,minimum width=0.39\textwidth,minimum height=18mm,
               align=center,fill=modsimgray},
  item/.style={draw,rounded corners,align=center,font=\small,minimum width=0.31\textwidth,
               minimum height=8mm},
  arrow/.style={-{Latex[length=2mm]},thick}
]
  \node[proc,fill=modsimblue!10,draw=modsimblue] (qt) {Qt parent process\\Runtime Inspector};
  \node[proc,fill=modsimcyan!10,draw=modsimcyan!80!black,right=of qt] (host)
    {MuJoCo child process\\authoritative runtime + native viewer};
  \node[item,below=of qt] (present) {Presenter, stable graph layout, event table};
  \node[item,below=of host] (owner) {RuntimeSession, scenario, ModelViewFactory};
  \node[item,below=of owner] (engine) {MjModel / MjData / passive viewer};
  \draw[arrow] (host) -- node[above,font=\scriptsize]{versioned immutable frames} (qt);
  \draw[arrow] (qt) -- node[below,font=\scriptsize]{initialize / stop} (host);
  \draw[arrow] (qt) -- (present);
  \draw[arrow] (host) -- (owner);
  \draw[arrow] (owner) -- (engine);
\end{tikzpicture}
\caption{Dual-window Runtime Inspector architecture. The protocol carries values only;
the Qt process never traverses mutable world or engine objects.}
\label{fig:runtime-processes}
\end{figure}

\subsection{Versioned local protocol}

The child emits newline-delimited strict JSON records prefixed by
\code{MODSIM\_RUNTIME/1}. Message kinds are hello, initialize, stop, status, frame,
error, and finished. The framer handles fragmented reads and rejects bad prefixes,
UTF-8, multiline documents, oversized records, and unterminated input. Standard output
is reserved for protocol data; standard error carries diagnostics that the parent
mirrors into its session log.

Shutdown is cooperative: the parent sends Stop, waits, then escalates to terminate and
kill only after deadlines. End-of-file also stops the child. The child exits the viewer,
shuts down the backend, sends a terminal record, and joins its input watcher. These
details prevent orphan viewers and interpreter aborts from blocked standard-input
threads.

\subsection{Presentation behavior}

The Qt-free presenter owns layout and selection, not runtime state. It assigns stable
deterministic node positions, keeps nodes in place across pose samples, renders parallel
edges as distinct curves, removes an edge selection when that edge disappears, and
retains node selection. It rejects regressing source stamps and gaps or conflicts in
event deltas.

The window shows graph, event table, scenario phase/detail, backend, simulation time,
module/assembly/connection counts, dock/undock outcomes, event count, and source
revisions. A target-speed label explicitly states the requested pacing factor. The
native viewer renders physical geometry and contact; the semantic window does not.

\subsection{Current visualization limits}

Only the module-topology graph is rendered. There are no frame, port, matrix, lattice,
or hybrid-docking panels; no time-series plots; no pause, restart, or scrub; no manual
joint controls; and no general scene editor. The Runtime Inspector remains a separate
window from Robot Pack Studio. These are presentation limitations rather than reasons
to weaken the immutable view boundary.

